\cotefile{15}
\scanpage{218}
% The typescript numbers its displayed equations at the left.
\makeatletter
\tagsleft@true
\makeatother
% archive sheet number: 217

\mspage{1 - 47 -} \hfill n\textsuperscript{o} 413

16. \uline{Mesure de Haar d'une variété de groupe}.

Conservons les notations de la fin du n\textsuperscript{o} précédent. L'espace $\Omega_e^n(G)$ est alors de dimension $1$ et il en résulte que \uline{l'espace des formes différentielles invariantes à gauche et de degré} $n=\dim G$ \uline{est de dimension} $1$. On en obtient une \uline{base} en prenant une base $\omega_1,\ldots,\omega_n$ de $\Omega_e^1(G)$ et en prenant le produit extérieur des formes $I^g_{\omega_k}$. Avec les notations de (30), on obtient en prenant $\omega_k=\mathrm{d}x_k$, l'expression d'une forme invariante à gauche $I^g_\omega$ de degré maximum au voisinage de $e$ :

\[
I^g_\omega(x)=(\text{dét}\,\partial m_k/\partial y_j(x^{-1};x))\,\mathrm{d}x_1\,\ldots\,\mathrm{d}x_n. \tag*{(31)}
\]

D'autre part, il est clair que l'espace des formes différentielles invariantes à gauche est invariant par translations à droite. Il existe donc une fonction de classe $C^{r-1}$ sur $G$, soit $\delta$, telle que :

\[
I^g_\omega(xy)=\delta(y)\,I^g_\omega\qquad (x,y\in G\,,\;\omega\in\Omega_e^n(G)). \tag*{(32)}
\]

Mais on a $\delta(y)\,I^g_\omega=I^g_\omega(xy)=I^g_\omega(y^{-1}xy)=I^g_{\omega'}$ avec $\omega'=(\operatorname{ad}y^{-1})_*(\omega)$. D'où :

\[
\delta(y)=(\text{dét}\,\operatorname{ad}y)^{-1}. \tag*{(33)}
\]

Supposons maintenant que $K$ soit \uline{localement compact} et soit $\omega\in\Omega_e^n(G)$. Il est clair que la mesure $\mu=|I^g_\omega|$ est une \uline{mesure de Haar à gauche}\unsure{.} sur le groupe localement compact $G$. Avec les notations de (30) et (31), on a donc pour $x$ voisin de $e$ :

\[
\mathrm{d}\mu(x)=\operatorname{mod}(\text{dét}\,\partial m_k/\partial y_j(x^{-1};x)))\,\mathrm{d}x_1\,\ldots\,\mathrm{d}x_n \tag*{(34)}
\]

où $\operatorname{mod}$ est le \uline{module} du corps localement compact $K$ (Int. VII, \S\ 1, n\textsuperscript{o} 10 et F. R. V.). D'autre part, (32) et (33) montrent que le \uline{module} du groupe localement compact $G$ (Int. VII, \S\ 1, n\textsuperscript{o} 3, déf. 3) n'est autre que

\[
\Delta(y)=\operatorname{mod}(\delta(y))=\operatorname{mod}(\text{dét}\,\operatorname{ad}y)^{-1}. \tag*{(35)}
\]
